\begin{center}
    {\fontsize{14}{16}\selectfont\bfseries ABSTRAK}
\end{center}

\vspace{0.5cm}

\begin{singlespace}
\noindent
Pencatatan transaksi pada operasional laundry yang masih bersifat manual seringkali menimbulkan risiko ketidakakuratan data dan kurangnya transparansi pembayaran. Penelitian ini bertujuan untuk mengembangkan Sistem \textit{Point of Sale} (POS) WellClean Laundry berbasis web yang mengintegrasikan fitur manajemen pelanggan, layanan, dan transaksi otomatis. Sistem ini dibangun menggunakan bahasa pemrograman Go dengan arsitektur yang mendukung efisiensi pemrosesan data. Salah satu fitur unggulan yang diimplementasikan adalah integrasi pembayaran digital melalui QRIS yang didukung oleh mekanisme \textit{webhook} untuk mendeteksi dana masuk secara \textit{real-time} dan menutup jendela transaksi secara otomatis. Guna menjamin keandalan sistem, dilakukan audit pengujian unit (\textit{unit testing}) menggunakan \textit{tool go tool cover}. Hasil pengujian menunjukkan fungsionalitas krusial seperti GopayHandler dan HashPassword mencapai tingkat validitas sebesar 100.0\%, dengan rata-rata keseluruhan \textit{statements} sistem sebesar 15.3\%. Implementasi ini terbukti dapat meningkatkan akurasi operasional dan memberikan kemudahan bagi kasir dalam memproses transaksi digital.

\vspace{0.8cm} % Jarak ke kata kunci

\noindent
\textbf{Kata Kunci:} \textit{Point of Sale, Go, QRIS, Unit Testing, Code Coverage.}
\end{singlespace}